\documentclass[a4paper,10pt]{article}

\usepackage[utf8]{inputenc}
\usepackage[T1]{fontenc}
\usepackage[ngerman]{babel}
\usepackage{amsmath}
\usepackage{amssymb}
\usepackage{xcolor}
\usepackage{enumitem}
\usepackage{geometry}
\geometry{a4paper,margin=1.9cm}

% ----- Light Mode -----
\definecolor{accent}{HTML}{1A4E8A}
\definecolor{accent2}{HTML}{8A5A00}
\definecolor{accent3}{HTML}{2E6E3A}
\definecolor{accent4}{HTML}{7A3E9A}
\definecolor{rulecol}{HTML}{B0B0B0}

\newcommand{\kapitel}[1]{%
  \par\vspace{7pt}
  {\large\bfseries\color{accent} #1}\par
  \noindent\textcolor{accent}{\rule{\linewidth}{0.8pt}}\par\vspace{2pt}}

\newcommand{\aufg}[2]{%
  \par\vspace{3pt}
  {\bfseries\color{accent2} #1}~~{\footnotesize\color{accent2}[#2]}\par\vspace{1pt}}

\newcommand{\teaser}[1]{%
  \par\vspace{2pt}{\color{accent4}\textbf{Denkanstoß.}}~#1\par}

\newcommand{\reflexion}{%
  \par\vspace{2pt}{\bfseries\color{accent3} Selbstreflexionsfragen}\par\vspace{1pt}}

\setlist[enumerate]{leftmargin=1.8em,itemsep=1pt,topsep=1pt,parsep=0pt}
\setlist[itemize]{leftmargin=1.3em,itemsep=1pt,topsep=1pt,parsep=0pt,label=\textcolor{accent3}{\textbullet}}
\setlength{\parindent}{0pt}
\setlength{\parskip}{3pt}

\begin{document}

\begin{center}
  {\Large\bfseries\color{accent} Numerische Methoden -- Übungsaufgaben}\\[1pt]
  {\small Vollständig aus den Vorlesungsaufschrieben von Prof.\ Dr.-Ing.\ Mark
   Schutera extrahiert \;·\; ins Deutsche übersetzt}
\end{center}
\noindent\textcolor{rulecol}{\rule{\linewidth}{0.4pt}}

{\footnotesize Diese Sammlung enthält \emph{alle} im Skript gestellten Aufgaben
-- benannte Übungen, Handrechnungen, Code-/Maschinen-Aufgaben, Margin-Aufgaben,
Denkanstöße und Selbstreflexionsfragen -- über alle sechs Kapitel. Ausgearbeitete
Musterlösungen sind \emph{bewusst nicht} enthalten. Lösungswege: siehe separates
Dokument \texttt{schutera\_loesungen}.}

% =====================================================================
\kapitel{Kapitel 1 \;--\; Einstieg in die Numerischen Methoden}

\aufg{1.1 \;--\; Größeres Gleichungssystem}{Handrechnung}
Mit der Größe des Systems steigt der Rechenaufwand. Lösen Sie nun dieses größere
System aus drei Gleichungen mit drei Unbekannten von Hand:
\[
  \begin{pmatrix}2&1&3\\1&4&2\\3&2&5\end{pmatrix}
  \begin{pmatrix}w_1\\w_2\\w_3\end{pmatrix}=\begin{pmatrix}14\\20\\32\end{pmatrix}
\]

\aufg{1.2 \;--\; Grid-Search}{Handrechnung}
Führen Sie eine Grid-Search für $w$ und $b$ mit Schrittweite $2{,}5$ über den
Bereich $w,b\in\{0;\,2{,}5;\,5;\,7{,}5;\,10\}$ durch. Berechnen Sie für jede
Kombination den Fehler $\mathrm{error}_1=|2w+b-11|$, $\mathrm{error}_2=|5w+b-13|$.
Werten Sie alle Kombinationen aus und wählen Sie das Paar $(w,b)$ mit dem
kleinsten akkumulierten Fehler.

\aufg{1.3 \;--\; Enter the machine}{Code}
Viele numerische Verfahren sind für die Umsetzung am Computer gedacht. Sie
erhalten ein vorgehostetes Jupyter-Notebook mit einer Python-Vorlage für diese
Übung. Nutzen Sie die Selbstreflexionsfragen unten, um sich beim Erkunden und
Experimentieren mit der Implementierung leiten zu lassen.

\teaser{Fällt Ihnen eine einfache Möglichkeit ein, das Verhältnis von Genauigkeit
zu Rechenaufwand unseres Grid-Search-Beispiels von oben zu verbessern?}

\reflexion
\begin{itemize}
  \item Warum ist die Wahl der Schrittweite $h$ beim Diskretisieren einer stetigen
        Funktion wichtig, und wie beeinflusst sie Genauigkeit und Rechenzeit der
        numerischen Lösung?
  \item Wie beeinflusst die Wahl des Intervalls $[a,b]$ die Ergebnisse der
        Diskretisierung und die Lage der numerisch gefundenen Extrema?
  \item Was sind die wesentlichen Unterschiede zwischen einer stetigen Funktion und
        ihrer diskretisierten Version, und welche Implikationen ergeben sich für das
        numerische Lösen mathematischer Probleme?
\end{itemize}

% =====================================================================
\kapitel{Kapitel 2 \;--\; Gleitkomma-Arithmetik}

\aufg{2.1 \;--\; Rundungsfehler}{Handrechnung}
Bekommen wir an einem einfachen Beispiel ein Gefühl für den Rundungsfehler.
Betrachten Sie die Division von $10$ durch $3$. Schreiben Sie die Dezimalentwicklung
aus.

\aufg{2.2 \;--\; Maschinenepsilon (hands on)}{Code}
Bevor Sie an den Rechner gehen: Überlegen Sie, wie Sie das Maschinenepsilon eines
beliebigen Systems für ein bestimmtes Gleitkommaformat per Programm bestimmen
würden. Rufen Sie sich die Definition des Maschinenepsilons in Erinnerung. Schreiben
Sie Ihre Überlegungen als Pseudocode auf. Was würde ein solches Programm auf einem
Festkomma- gegenüber einem Gleitkommasystem zeigen? Schreiben Sie Ihre Gedanken auf.
Probieren Sie dann verschiedene Datentypen am Rechner im Code aus. Notieren Sie Ihre
Beobachtungen und reflektieren Sie sie. Wenn Sie einen Taschenrechner bauen würden --
welchen Datentyp würden Sie wählen und warum?

\aufg{2.3 \;--\; Auslöschung (loss of significance)}{Code}
Überlegen Sie nun, wie Sie Auslöschung an Ihrem Rechner demonstrieren würden.
Schreiben Sie Ihre Überlegungen als Pseudocode auf, bevor Sie weiterlesen.

\aufg{2.4 \;--\; Argumentation}{Konzept}
Überlegen Sie, wie das Berechnen von $f(\epsilon)=\dfrac{1}{a+\epsilon-b}$ für
kleines $\epsilon$ und $a=b$ dies leisten würde. Was erwarten Sie zu sehen, wenn
$\epsilon$ sehr klein ist?
\par\textit{Margin:} Was passiert für $a>b$? Denken Sie in signifikanten Stellen.

\teaser{Fallen Ihnen Metriken für numerische Verfahren ein, basierend auf den
Approximationen und Fehlern, die wir besprochen haben?}

\reflexion
\begin{itemize}
  \item Wie ist die Gleitkommadarstellung aufgebaut, und was sind ihre Komponenten?
  \item Was ist das Maschinenepsilon, und wie hängt es mit der numerischen
        Genauigkeit zusammen?
  \item Wie wirken sich diese Konzepte in der Praxis auf numerische Berechnungen aus?
\end{itemize}

% =====================================================================
\kapitel{Kapitel 3 \;--\; Fehleranalyse}

\aufg{3.1 \;--\; Minimum der Sinusfunktion}{Handrechnung}
Bestimmen Sie das Minimum der Sinusfunktion auf dem Intervall $[-3,1]$ noch einmal.
Diskretisieren Sie die Funktion erneut mit zwei verschiedenen Schrittweiten $h=1$ und
$h=0{,}2$. Nehmen Sie weiter an, dass die Eingabedaten $x$ durch Messfehler um
$\tilde x=x\pm0{,}25$ gestört sind und eine Genauigkeit von zwei Dezimalstellen
besitzen. Analysieren Sie nun Kondition, Stabilität, Konsistenz und Konvergenz des zur
Minimumsuche verwendeten numerischen Verfahrens in beiden Fällen. Quantifizieren Sie
die Eigenschaften mithilfe der im vorigen Abschnitt angegebenen Formeln.

\aufg{3.2 \;--\; Stabilität und Kondition}{Beweis}
\textit{Margin:} Die Stabilität entspricht der Kondition des Systems für $h\to0$.
Zeigen Sie das.

\aufg{3.3 \;--\; Compute-getriebene Fehleranalyse (hands on)}{Code}
Brute-forcen wir uns durch die Fehleranalyse beim Lösen des Zwei-Gleichungs-Systems
aus Kapitel 1. Sie bestimmen und optimieren Kondition, Stabilität, Konsistenz und
Konvergenz der numerischen Lösung, indem Sie das numerische Verfahren optimieren.

\reflexion
\begin{itemize}
  \item Wie gut konditioniert ist das Zwei-Gleichungs-System verglichen mit der
        numerischen Lösung des Minimums der Sinusfunktion?
  \item Konvergiert die Stabilität beim Zwei-Gleichungs-System mit kleineren
        Schrittweiten? Wogegen?
  \item Wie wirken sich Störungen der Eingabedaten auf die numerische Lösung des
        Zwei-Gleichungs-Systems aus? Gibt es ein Zusammenspiel mit der Schrittweite?
  \item Vergleichen Sie Konvergenz mit Stabilität und Konsistenz. Was beobachten Sie?
        Können Sie Ihre Beobachtungen in Begriffen von Bias und Varianz ausdrücken?
  \item Beobachten Sie für immer kleinere Schrittweiten eine Grenze der Genauigkeit
        der numerischen Lösung? Wenn ja, warum?
  \item Welches weitere Problem beobachten Sie beim Verfeinern der Schrittweite?
\end{itemize}

\teaser{Bisher haben wir Brute-Force-Lösungen verwendet. Fällt Ihnen eine Möglichkeit
ein, Brute-Force-Methoden zu verfeinern, um mit weniger Aufwand bessere Ergebnisse zu
erzielen? Nutzen Sie den Code dieser Vorlesung, um Ihre Idee zu entwerfen und zu
testen.}

% =====================================================================
\kapitel{Kapitel 4 \;--\; Newton-Verfahren}

\aufg{4.0 \;--\; Denksportaufgabe: Versagensmodi}{Schaubild}
Das Newton-Verfahren kann auf verschiedene Weisen versagen. Als Denksportaufgabe ist
es eine schöne Übung, die Versagensmodi in diesem Schaubild zu visualisieren
(Abb.: $f(x)=x^3-x$).

\aufg{4.1 \;--\; Newton von Hand}{Handrechnung}
Finden Sie die Nullstelle von $f(x)=x^3-x=0$, indem Sie das Newton-Verfahren
geometrisch und anschließend von Hand anwenden. Die Nullstellen liegen bei
$x=-1,0,1$. Suchen wir die Nullstelle bei $x=1$, ausgehend von $x_0=1{,}4$.

\aufg{4.2 \;--\; Sekantenverfahren von Hand}{Handrechnung}
Wenden Sie das Sekantenverfahren auf dasselbe Problem mit $x_0=1{,}2$, $x_1=1{,}4$
an. Wieder zuerst geometrisch, dann von Hand.

\aufg{4.3 \;--\; Geometrische Fehlersuche}{Handrechnung}
Finden Sie Startpunkte $x_0$ für verschiedene Versagensmodi: einen, bei dem das
Newton-Verfahren wegen verschwindender Ableitung versagt, einen, bei dem es zu einer
Nicht-Ziel-Nullstelle konvergiert, und einen, bei dem es zwischen Punkten oszilliert.

\aufg{4.4 \;--\; Enter the machine}{Code}
Betrachten wir Konvergenz und Konvergenzraten von Grid-Search, Newton- und
Sekantenverfahren in Python. Experimentieren Sie dann mit verschiedenen Startpunkten
und beobachten Sie die Versagensmodi in der Praxis. Versuchen Sie, die Startpunkte
zuerst geometrisch zu finden, dann analytisch, bevor Sie sie schließlich im Code
verifizieren.

\reflexion
\begin{itemize}
  \item Was ist der Hauptvorteil des Newton-Verfahrens gegenüber der Grid-Search?
  \item Warum genügt eine lineare Approximation zum Finden von Nullstellen? Was sagt
        uns die Taylor-Entwicklung über diese Wahl?
  \item Warum ist quadratische Konvergenz so mächtig? Wie viele Iterationen bräuchte
        Newton, um von Fehler $10^{-1}$ auf $10^{-16}$ zu kommen?
  \item In welchen Situationen würden Sie das Sekantenverfahren dem Newton vorziehen?
        Warum nicht in anderen?
\end{itemize}

\teaser{Wie können wir sicherstellen, dass wir die globale Lösung finden und nicht
nur eine lokale?}

% =====================================================================
\kapitel{Kapitel 5 \;--\; Globale Optimierung}

\aufg{5.0 \;--\; Newton auf Shekel, weiterer Startpunkt}{Code}
Wenden wir das Newton-Verfahren auf unsere 1D-Shekel-Funktion von verschiedenen
Startpunkten an und beobachten, wohin das Verfahren konvergiert. Wir zeigen es für
einen Startpunkt; probieren Sie selbst einen weiteren, und dann diskutieren wir das
allgemeine Verhalten.
\par\textit{Margin:} Wie würden Sie die Update-Regel abändern, um stattdessen Maxima
zu finden?

\aufg{5.1 \;--\; Random Restarts}{Handrechnung}
Angenommen, Sie optimieren eine Funktion mit 5 lokalen Minima, und das Einzugsgebiet
des globalen Minimums deckt $15\%$ des Suchraums ab ($p=0{,}15$). Wie groß ist die
Wahrscheinlichkeit, das globale Minimum mit $K=10$ Random Restarts zu finden?
\par\textit{Hinweis:} Nutzen Sie $P=1-(1-p)^K$ und lösen Sie nach $K$ auf:
$K=\dfrac{\ln(1-P)}{\ln(1-p)}$.
\par\textbf{(Folge)} Wie viele Restarts brauchen wir für $99\%$? \quad Berechnen Sie
nun $K$ für $p=0{,}05$ und $p=0{,}01$ beim selben $99\%$-Ziel.

\aufg{5.2 \;--\; Eigene 1D-Shekel-Funktion}{Code}
An Ihrem Rechner entwerfen Sie zunächst Ihre eigene 1D-Shekel-Funktion. Halten Sie
sie anfangs einfach, machen Sie sie später komplexer. Machen Sie sich dann zuerst mit
den bisher gelernten lokalen Optimierungsverfahren vertraut und wenden Sie sie auf
Ihre Funktion an. Zeigen Sie, wo das Newton-Verfahren (erneut) versagt. Untersuchen
Sie dann verschiedene $x_0$ und sehen Sie, wie sich das Konvergenzverhalten ändert.
Denken Sie wieder über die Loss-Landschaft und Basin-Grenzen nach.

\aufg{5.3 \;--\; Die Suche zu Minima lenken}{Konzept + Code}
Erinnern Sie sich an den $|f''|$-Trick, der das Vorzeichen der Krümmung im Nenner
umkehrt und den Schritt zwingt, stets bergab zu gehen. Überlegen Sie, was das
geometrisch für die Tangentenkonstruktion bedeutet. Probieren Sie es dann im Code aus.

\aufg{5.4 \;--\; Einzugsgebiet (Basin of Attraction)}{Schaubild}
Betrachten Sie erneut die Plots der Shekel-Funktion, markieren Sie die Einzugsgebiete
jedes lokalen Minimums und schätzen Sie ihre relativen Größen ab. Welches ist das
globale Minimum? Wie hängt das mit der Wahrscheinlichkeit zusammen, es per Random
Restarts zu finden?

\aufg{5.5 \;--\; Lokalen Minima entkommen}{Code}
Setzen Sie Random Restarts und Basin Hopping auf Ihrer eigenen 1D-Shekel-Funktion ein
und vergleichen Sie, wie viele Funktionsiterationen jede Methode bis zum globalen
Minimum braucht. Wie anfällig sind sie, in lokalen Minima hängenzubleiben? Wie
beeinflusst die Wahl der Schrittweitenverteilung die Performance des Basin Hopping?

\aufg{5.6 \;--\; Nicht-konvexe Becken}{Konzept}
Shekel ist im Allgemeinen nicht-konvex, die einzelnen Becken jedoch schon; notieren
Sie eine Funktion, bei der die Einzugsgebiete nicht so leicht abzugrenzen sind.

\aufg{5.7 \;--\; Adaptive Schrittweite}{Konzept + Code}
Fällt Ihnen ein geschickter Weg ein, die Schrittweitenverteilung während der
Optimierung automatisch anzupassen? Was wäre eine gute Heuristik? Wie könnten Sie
katastrophale Sprünge verhindern? Implementieren Sie es.

\aufg{5.8 \;--\; Noise and landscape engineering}{Konzept}
Die Loss-Landschaft ist nicht statisch, sondern soll geformt und gestaltet werden.
Indem man die Loss-Landschaft auf Mini-Batches (unterschiedliche Auswahl von Punkten)
approximiert \dots

\aufg{5.9 \;--\; Code-Übung: Restarts vs.\ Basin Hopping}{Code}
Implementieren Sie Random Restarts und Basin Hopping auf der 1D-Shekel-Funktion.
Vergleichen Sie, wie viele Funktionsauswertungen jede Methode braucht, um das globale
Minimum bei $x\approx4$ zu finden.

\aufg{5.10 \;--\; Einzugsgebiete kartieren}{Handrechnung}
Kartieren wir schließlich die Einzugsgebiete. Betrachten Sie die 1D-Shekel-Funktion
und nehmen Sie an, ein lokaler Optimierer konvergiere stets zum nächstgelegenen
Foxhole.
\begin{enumerate}
  \item[(a)] Skizzieren Sie (grob) die Einzugsgebiete für die drei Foxholes bei
        $x=0,4,7$. Wo liegen die Basin-Grenzen?
  \item[(b)] Wenn Basin Hopping einen Sprung $\delta\sim\mathrm{Uniform}(-3,3)$
        ausgehend vom lokalen Minimum $x^\circ=7$ verwendet, wie groß ist die
        Wahrscheinlichkeit, dass ein einzelner Sprung im Becken des globalen Minimums
        landet?
  \item[(c)] Warum könnte Basin Hopping das globale Minimum hier schneller finden als
        Random Restarts?
\end{enumerate}

\reflexion
\begin{itemize}
  \item Warum versagt lokale Optimierung auf nicht-konvexen Landschaften wie Shekels
        Foxholes?
  \item Wie skaliert der erforderliche Aufwand mit der Größe des globalen Beckens bei
        Random Restarts?
  \item Wie unterscheidet sich Basin Hopping von Random Restarts?
  \item Was sagt uns $f''(x_n)$ über die Loss-Landschaft, was charakterisiert es, und
        wie können wir es nutzbar machen?
  \item Wie hilft das Approximieren der Loss-Landschaft mit verrauschten Schätzungen
        (Mini-Batches), lokalen Optima zu entkommen?
  \item Wählen Sie $x_0$ so, dass es im Becken eines lokalen Minimums liegt; finden Sie
        nun verschiedene Wege, mit denen Ihr Optimierungsverfahren entkommen kann --
        wie wirksam ist welcher?
\end{itemize}

\teaser{Was, wenn wir nun einer hochfrequenten Version von Shekels Foxholes
gegenüberstehen? Welches Problem tritt dabei auf?}

% =====================================================================
\kapitel{Kapitel 6 \;--\; Numerische Integration}

\aufg{6.1 \;--\; Exakter Wert}{Handrechnung}
Bleiben wir bei unserem Beispiel aus dem Methodenabschnitt und berechnen das Integral
von $\sin(x)$ von $-2$ bis $-1$ mit verschiedenen Methoden. Um den Fehler jeder Methode
quantifizieren zu können, brauchen wir den exakten Wert. Berechnen wir also den exakten
Wert des Integrals.

\aufg{6.2 \;--\; Mittelpunktsregel von Hand}{Handrechnung}
Berechnen Sie $\int_{-2}^{-1}\sin(x)\,dx$ mit der Mittelpunktsregel mit $n=1$ Intervall.

\aufg{6.3 \;--\; Mittelpunktsregel mit mehr Intervallen}{Handrechnung}
Berechnen Sie nun die Mittelpunktsregel mit $n=2$ Intervallen, d.\,h.\ wir haben zwei
Mittelpunkte bei $-1{,}75$ und $-1{,}25$.
\par\textit{Margin:} Weiter lohnt es sich, die Approximation an den Endpunkten des
Intervalls zu berechnen und den Fehler zu vergleichen.

\aufg{6.4 \;--\; Trapezregel von Hand}{Handrechnung}
Berechnen Sie mit der Trapezregel $\int_{-2}^{-1}\sin(x)\,dx$ mit $n=2$ Intervallen.
Die Stützstellen sind $-2,-1{,}5$ und $-1$.

\aufg{6.5 \;--\; Simpson-Regel von Hand}{Handrechnung}
Berechnen Sie mit der Simpson-Regel $\int_{-2}^{-1}\sin(x)\,dx$ mit $n=2$ Intervallen.
Die Stützstellen sind $-2,-1{,}5$ und $-1$.
\par\textit{Margin:} Probieren Sie mehr Intervalle, sehen Sie, wie der Fehler abnimmt,
und prüfen Sie, ob Sie sich mit zusammengesetzten Simpson- und Trapezregeln auskennen.

\aufg{6.6 \;--\; Monte Carlo von Hand}{Handrechnung}
Schätzen Sie $\int_{-2}^{-1}\sin(x)\,dx$ per Monte Carlo mit $N=4$ Zufallsstichproben,
was etwa dem Rechenaufwand der Trapezmethode entspricht. Nehmen Sie an, die
Zufallsstichproben sind: $X_1=-1{,}95$, $X_2=-1{,}55$, $X_3=-1{,}15$, $X_4=-1{,}05$.

\aufg{6.7 \;--\; Monte Carlo: $N$ erhöhen}{Konzept}
Was passiert, wenn wir $N$ auf 16 erhöhen? \dots\ Erläutern Sie kurz, wie der Fehler
des Monte-Carlo-Schätzers dimensionsunabhängig ist.

\aufg{6.8 \;--\; Simpson-Gewichte via Lagrange}{Herleitung}
\textit{Margin:} Probieren Sie es mit ein paar Punkten. Beginnen Sie mit $x=0$ und
sehen Sie, wie die Linearfaktoren ausfallen. (Im Skript wird die Rechnung übersprungen:
Leiten Sie die Quadraturgewichte durch Integration der Lagrange-Basispolynome über
$[a,b]$ selbst her.)
\par\textit{Margin:} Was passiert, wenn Sie eine lineare Funktion mit einer
quadratischen approximieren?
\par\textit{Margin:} Was, wenn wir über ein Intervall integrieren, das zu destruktiver
Interferenz der zugrundeliegenden höherfrequenten Sinuskurve führt? Denken Sie an
Wellenlängen.

\aufg{6.9 \;--\; Höhere Dimensionen (on your machine)}{Code}
Implementieren Sie Monte-Carlo-Integration für ein $d$-dimensionales Integral, etwa die
Integration einer multivariaten Gauß-Funktion über einen Hyperwürfel. Während Sie $d$
langsam erhöhen, beobachten Sie, wie sich der Fehler für verschiedene Methoden verhält.
Beobachten Sie zudem, wie gitterbasierte Quadraturmethoden mit wachsendem $d$
rechnerisch undurchführbar werden. Optional: Sehen Sie, wie ähnliche Effekte in der
Optimierung auftreten, indem Sie Gradientenabstieg und Mini-Batch-SGD auf einer
hochdimensionalen Funktion implementieren. Achten Sie nur auf die Rechenzeiten.

\reflexion
\begin{itemize}
  \item Wie vergleichen sich die Fehler von Mittelpunkts-, Trapez-, Simpson-Regel und
        Monte Carlo?
  \item Warum hat die Mittelpunktsregel eine bessere Genauigkeit als die Links- oder
        Rechts-Endpunkt-Regel?
  \item Wie verbessern zusammengesetzte Methoden die Genauigkeit, und was ist der
        Trade-off?
  \item Warum versagen deterministische Quadraturmethoden in hohen Dimensionen?
  \item Inwiefern ist die Mittelpunktsregel ein Spezialfall der Monte-Carlo-Schätzung?
  \item Was ist die Intuition dahinter, dass Monte Carlo in hohen Dimensionen nicht
        ``explodiert''?
  \item Wann würden Sie die Simpson-Regel gegenüber Monte Carlo verwenden?
  \item Wie hängt der Mini-Batch-Mittelwert in SGD mit der Monte-Carlo-Schätzung
        zusammen?
  \item Wie hilft die Gradientenschätzung in SGD, lokalen Minima zu entkommen?
\end{itemize}

\teaser{Warum werden hochgradige Quadraturmethoden instabil, wenn sie Polynome
niedrigen Grades approximieren?}

\end{document}
